Since (112.7) does not involve $p_x$, whose values form a continuum, every energy level has a continuous degeneracy. This degeneracy becomes finite, however, when the motion in the $xy$ plane is confined to a large but finite area $S=L_xL_y$. The number of distinct (now discrete) values of $p_x$ in an interval $\Delta p_x$ is $(L_x/2\pi\hbar)\Delta p_x$. All $p_x$ for which the orbit centre lies within $S$ are allowed (the orbit radius is neglected in comparison with the large length $L_y$). The conditions $0<y_0<L_y$ give $\Delta p_x=eHL_y/c$. Consequently, for fixed $n$ and $p_z$, the number of states is $eHS/2\pi\hbar c$. If the region of motion is also bounded along the $z$ axis, with length $L_z$, then an interval $\Delta p_z$ contains $(L_z/2\pi\hbar)\Delta p_z$ possible values of $p_z$, and the number of states in that interval is

